% GENERATED FILE. Edit canonical lessons, then run npm run generate:books.
% canonical-source-hash: fnv1a64:f29c32f03b7e402f
% canonical-lessons: JA-C365-keizai, JA-C365-keiki, JA-C365-toushi, JA-C365-nenkin, JA-C365-heikin

\chapter{The Economy}
\label{ch:ja-the-economy}
\hlchaptermodality{365}

\begin{chapteropening}
\textbf{By the end of this chapter:} \emph{I can say what would happen if the economy got better.}

Complete the last lesson of chapter 365: I can say what would happen if the economy got better.
\end{chapteropening}

\section[\texorpdfstring{\ja{けいざい} (keizai)}{けいざい (keizai)}]{\ja{けいざい} (keizai) — the economy}

\label{lesson:JA-C365-keizai}



\begin{quote}
Before the new one: say the Japanese for to be in charge of, then the Japanese for responsibility.
\end{quote}

\subsection*{You'll want to know: \ja{けいざい}}
\textbf{\ja{けいざい}} — \emph{keizai} — \textquotedblleft{}the economy\textquotedblright{}.

\begin{quote}
\textbf{\ja{けいざいが} \ja{よく} \ja{なったら、しごとが} \ja{ふえるでしょう。}} — \emph{keizai ga yoku nattara, shigoto ga fueru deshō} — \textquotedblleft{}If the economy gets better, there will probably be more jobs.\textquotedblright{}
\end{quote}

\begin{cousinweb}
\textbf{\ja{けいざい}} (\emph{keizai}) is Sino-Japanese (kango): two Chinese characters read with their on'yomi, readings borrowed from Chinese. They are \emph{kei}, \textquotedblleft{}pass thru\textquotedblright{} (Middle Chinese \emph{keng}) and \emph{zai}, \textquotedblleft{}settle\textquotedblright{} (\emph{tsejH}) (KANJIDIC2; Middle Chinese in Baxter's transcription).
\end{cousinweb}

\begin{grammarlens}[title={What you've built}]
One more word for saying what happens if.
\end{grammarlens}

\subsection*{Your turn}
\begin{itemize}
\raggedright
  \item \emph{Say it:} \emph{keizai}
  \item \emph{Say it:} \emph{keizai}, once more
  \item \emph{Say it:} \emph{keizai ga yoku nattara, shigoto ga fueru deshō} — If the economy gets better, there will probably be more jobs.
  \item \emph{Recall:} say the Japanese for to be in charge of, then the Japanese for responsibility, then say \emph{keizai} again
\end{itemize}

\begin{tcolorbox}[breakable,colback=teal!4,colframe=teal!35!black,title={Before you move on}]
What does \ja{けいざい} mean? (\textquotedblleft{}The economy\textquotedblright{}.) Say it once more.
\end{tcolorbox}
\section[\texorpdfstring{\ja{けいき} (keiki)}{けいき (keiki)}]{\ja{けいき} (keiki) — business conditions}

\label{lesson:JA-C365-keiki}



\begin{quote}
Before the new one: say the Japanese for responsibility, then the Japanese for the economy.
\end{quote}

\subsection*{You'll want to know: \ja{けいき}}
\textbf{\ja{けいき}} — \emph{keiki} — \textquotedblleft{}business conditions\textquotedblright{}.

\begin{quote}
\textbf{\ja{けいきが} \ja{よく} \ja{なったら、みせが} \ja{ふえるでしょう。}} — \emph{keiki ga yoku nattara, mise ga fueru deshō} — \textquotedblleft{}If business picks up, there will probably be more shops.\textquotedblright{}
\end{quote}

\begin{cousinweb}
\textbf{\ja{けいき}} (\emph{keiki}) is Sino-Japanese (kango): two Chinese characters read with their on'yomi, readings borrowed from Chinese. They are \emph{kei}, \textquotedblleft{}scenery\textquotedblright{} and \emph{ki}, \textquotedblleft{}atmosphere\textquotedblright{} (Middle Chinese \emph{khj+jH}) (KANJIDIC2; Middle Chinese in Baxter's transcription).
\end{cousinweb}

\begin{grammarlens}[title={What you've built}]
One more word for saying what happens if.
\end{grammarlens}

\subsection*{Your turn}
\begin{itemize}
\raggedright
  \item \emph{Say it:} \emph{keiki}
  \item \emph{Say it:} \emph{keiki}, once more
  \item \emph{Say it:} \emph{keiki ga yoku nattara, mise ga fueru deshō} — If business picks up, there will probably be more shops.
  \item \emph{Recall:} say the Japanese for responsibility, then the Japanese for the economy, then say \emph{keiki} again
\end{itemize}

\begin{tcolorbox}[breakable,colback=teal!4,colframe=teal!35!black,title={Before you move on}]
What does \ja{けいき} mean? (\textquotedblleft{}Business conditions\textquotedblright{}.) Say it once more.
\end{tcolorbox}
\section[\texorpdfstring{\ja{とうし} (tōshi)}{とうし (tōshi)}]{\ja{とうし} (tōshi) — an investment}

\label{lesson:JA-C365-toushi}



\begin{quote}
Before the new one: say the Japanese for the economy, then the Japanese for business conditions.
\end{quote}

\subsection*{You'll want to know: \ja{とうし}}
\textbf{\ja{とうし}} — \emph{tōshi} — \textquotedblleft{}an investment\textquotedblright{}.

\begin{quote}
\textbf{\ja{とうしの} \ja{けっかは、けいきに} \ja{よります。}} — \emph{tōshi no kekka wa, keiki ni yorimasu} — \textquotedblleft{}How an investment turns out depends on business conditions.\textquotedblright{}
\end{quote}

\begin{cousinweb}
\textbf{\ja{とうし}} (\emph{tōshi}) is Sino-Japanese (kango): two Chinese characters read with their on'yomi, readings borrowed from Chinese. They are \emph{tō}, \textquotedblleft{}throw\textquotedblright{} (Middle Chinese \emph{duw}) and \emph{shi}, \textquotedblleft{}resources\textquotedblright{} (\emph{tsij}) (KANJIDIC2; Middle Chinese in Baxter's transcription).
\end{cousinweb}

\begin{grammarlens}[title={What you've built}]
One more word for saying what happens if.
\end{grammarlens}

\subsection*{Your turn}
\begin{itemize}
\raggedright
  \item \emph{Say it:} \emph{tōshi}
  \item \emph{Say it:} \emph{tōshi}, once more
  \item \emph{Say it:} \emph{tōshi no kekka wa, keiki ni yorimasu} — How an investment turns out depends on business conditions.
  \item \emph{Recall:} say the Japanese for the economy, then the Japanese for business conditions, then say \emph{tōshi} again
\end{itemize}

\begin{tcolorbox}[breakable,colback=teal!4,colframe=teal!35!black,title={Before you move on}]
What does \ja{とうし} mean? (\textquotedblleft{}An investment\textquotedblright{}.) Say it once more.
\end{tcolorbox}
\section[\texorpdfstring{\ja{ねんきん} (nenkin)}{ねんきん (nenkin)}]{\ja{ねんきん} (nenkin) — a pension}

\label{lesson:JA-C365-nenkin}



\begin{quote}
Before the new one: say the Japanese for business conditions, then the Japanese for an investment.
\end{quote}

\subsection*{You'll want to know: \ja{ねんきん}}
\textbf{\ja{ねんきん}} — \emph{nenkin} — \textquotedblleft{}a pension\textquotedblright{}.

\begin{quote}
\textbf{\ja{しごとを} \ja{やめたら、ねんきんで} \ja{くらします。}} — \emph{shigoto o yametara, nenkin de kurashimasu} — \textquotedblleft{}When I stop working, I will live on my pension.\textquotedblright{}
\end{quote}

\begin{cousinweb}
\textbf{\ja{ねんきん}} (\emph{nenkin}) is Sino-Japanese (kango): two Chinese characters read with their on'yomi, readings borrowed from Chinese. They are \emph{nen}, \textquotedblleft{}year\textquotedblright{} (Middle Chinese \emph{nen}) and \emph{kin}, \textquotedblleft{}gold\textquotedblright{} (\emph{kim}) (KANJIDIC2; Middle Chinese in Baxter's transcription).
\end{cousinweb}

\begin{grammarlens}[title={What you've built}]
One more word for saying what happens if.
\end{grammarlens}

\subsection*{Your turn}
\begin{itemize}
\raggedright
  \item \emph{Say it:} \emph{nenkin}
  \item \emph{Say it:} \emph{nenkin}, once more
  \item \emph{Say it:} \emph{shigoto o yametara, nenkin de kurashimasu} — When I stop working, I will live on my pension.
  \item \emph{Recall:} say the Japanese for business conditions, then the Japanese for an investment, then say \emph{nenkin} again
\end{itemize}

\begin{tcolorbox}[breakable,colback=teal!4,colframe=teal!35!black,title={Before you move on}]
What does \ja{ねんきん} mean? (\textquotedblleft{}A pension\textquotedblright{}.) Say it once more.
\end{tcolorbox}
\section[\texorpdfstring{\ja{へいきん} (heikin)}{へいきん (heikin)}]{\ja{へいきん} (heikin) — an average}

\label{lesson:JA-C365-heikin}



\begin{quote}
Before the new one: say the Japanese for an investment, then the Japanese for a pension.
\end{quote}

\subsection*{You'll want to know: \ja{へいきん}}
\textbf{\ja{へいきん}} — \emph{heikin} — \textquotedblleft{}an average\textquotedblright{}.

\begin{quote}
\textbf{\ja{へいきんを} \ja{けいさんしたら、おしえて} \ja{ください。}} — \emph{heikin o keisan shitara, oshiete kudasai} — \textquotedblleft{}When you have worked out the average, please tell me.\textquotedblright{}
\end{quote}

\begin{cousinweb}
\textbf{\ja{へいきん}} (\emph{heikin}) is Sino-Japanese (kango): two Chinese characters read with their on'yomi, readings borrowed from Chinese. They are \emph{hei}, \textquotedblleft{}even\textquotedblright{} (Middle Chinese \emph{bjaeng}) and \emph{kin}, \textquotedblleft{}level\textquotedblright{} (\emph{kjwin}) (KANJIDIC2; Middle Chinese in Baxter's transcription).
\end{cousinweb}

\begin{grammarlens}[title={What you've built}]
One more word for saying what happens if.
\end{grammarlens}

\subsection*{Your turn}
\begin{itemize}
\raggedright
  \item \emph{Say it:} \emph{heikin}
  \item \emph{Say it:} \emph{heikin}, once more
  \item \emph{Say it:} \emph{heikin o keisan shitara, oshiete kudasai} — When you have worked out the average, please tell me.
  \item \emph{Recall:} say the Japanese for an investment, then the Japanese for a pension, then say \emph{heikin} again
\end{itemize}

\begin{tcolorbox}[breakable,colback=teal!4,colframe=teal!35!black,title={Before you move on}]
What does \ja{へいきん} mean? (\textquotedblleft{}An average\textquotedblright{}.) Say it once more.
\end{tcolorbox}
